\documentclass[aspectratio=169]{beamer}
\usepackage{iftex}
\ifPDFTeX
  \usepackage[T1]{fontenc}
  \usepackage{lmodern}
\fi
\usepackage[czech]{babel}
\usepackage{listings}
\usepackage{tikz}
\usetikzlibrary{arrows.meta,positioning,babel}

\newcommand{\CVUTTemplatePath}{../../../utils/presentation-template}
\usepackage{\CVUTTemplatePath/beamerthemeCVUT}
\setbeamertemplate{footline}[frame number]
\beamertemplatenavigationsymbolsempty
\newcommand{\LANG}{CZ}
\newcommand{\FRONTPAGE}{dark}
\title{Ukládání a serializace dat}
\subtitle{Základy programování\\Matouš Cejnek\\U12110, FS, ČVUT v Praze}
\author{} \institute{} \date{}
\renewcommand{\headlineA}{} \renewcommand{\headlineB}{} \renewcommand{\headlineC}{}

\definecolor{CodeBlue}{RGB}{38,83,125}
\definecolor{SoftGray}{RGB}{242,244,247}
\definecolor{CodeGreen}{RGB}{38,125,94}
\definecolor{WarningRed}{RGB}{170,45,45}
\setbeamerfont{normal text}{size=\large}
\setbeamerfont{frametitle}{size=\LARGE}
\setbeamercolor{block title}{fg=white,bg=CodeBlue}
\setbeamercolor{block body}{fg=black,bg=SoftGray}
\newcommand{\term}[1]{\textcolor{CodeBlue}{\textbf{#1}}}
\lstdefinestyle{python}{language=Python,basicstyle=\ttfamily\large,
  keywordstyle=\color{CodeBlue}\bfseries,stringstyle=\color{CodeGreen},
  commentstyle=\color{gray},showstringspaces=false,columns=fullflexible,keepspaces=true}
\lstdefinestyle{pythonSmall}{style=python,basicstyle=\ttfamily\normalsize}
\lstdefinestyle{data}{basicstyle=\ttfamily\normalsize,showstringspaces=false,
  columns=fullflexible,keepspaces=true}

\begin{document}

% 0:00--0:01
\begin{frame}
  \clearpage\thispagestyle{empty}\titlepage
\end{frame}

% 0:01--0:04
\begin{frame}{Co si dnes odnesete?}
  \begin{itemize}
    \item Vysvětlit serializaci a deserializaci.
    \item Rozlišit datový model, formát a soubor.
    \item Uložit a načíst data pomocí modulů \texttt{json} a \texttt{csv}.
    \item Předpovědět, které datové typy se při převodu zachovají.
    \item Zvolit formát podle tvaru dat, příjemce a důvěryhodnosti zdroje.
    \item Bezpečně posoudit použití binárního formátu \texttt{pickle}.
  \end{itemize}
\end{frame}

\begin{frame}{Na co navazujeme?}
  \begin{itemize}
    \item Slovníky a seznamy nesou strukturovaná data v paměti.
    \item \texttt{with open(...)} bezpečně spravuje otevřený soubor.
    \item UTF-8 převádí text na bajty a zpět.
    \item Moduly standardní knihovny importujeme a používáme přes jejich API.
    \item Očekávatelné chyby umíme zachytit pomocí \texttt{try}/\texttt{except}.
  \end{itemize}
  \begin{block}{Z domácí přípravy}
    Proč se po načtení knih z CSV změnily \texttt{year} a \texttt{available} na řetězce?
  \end{block}
\end{frame}

\begin{frame}{Objekt v paměti po ukončení programu zmizí}
  \begin{itemize}
    \item Python pracuje s objekty: seznamy, slovníky, čísly a instancemi tříd.
    \item Soubor nebo síť přenáší pouze posloupnost znaků či bajtů.
    \item Příští spuštění programu potřebuje strukturu z této posloupnosti obnovit.
  \end{itemize}
  \begin{block}{Otázka}
    Jak zapíšeme strukturu dat tak, aby její význam znal i budoucí program?
  \end{block}
\end{frame}

% 0:04--0:08
\begin{frame}{Serializace je převod reprezentace}
  \centering
  \begin{tikzpicture}[>=Latex,node distance=18mm,font=\normalsize,
      box/.style={draw,rounded corners,fill=SoftGray,minimum width=28mm,minimum height=12mm,align=center}]
    \node[box] (object) {Python objekty\\\texttt{list}, \texttt{dict}, \texttt{int}};
    \node[box,right=of object] (encoded) {text nebo bajty\\podle formátu};
    \node[box,right=of encoded] (restored) {nové Python\\objekty};
    \draw[->,thick,CodeBlue] (object)--node[above,font=\small]{serializace}(encoded);
    \draw[->,thick,CodeBlue] (encoded)--node[above,font=\small]{deserializace}(restored);
  \end{tikzpicture}
  \vspace{7mm}

  \term{Deserializace nevzkřísí původní objekt.} Vytvoří nové objekty z uložené reprezentace.
\end{frame}

\begin{frame}{Formát není totéž co místo uložení}
  \begin{description}
    \item[Datový model] Jaké hodnoty a vztahy potřebujeme vyjádřit?
    \item[Formát] Jak budou data zapsána? JSON, CSV, binární formát.
    \item[Médium] Kde reprezentace leží? Soubor, databáze, síťová zpráva.
    \item[Kódování] Jak se text převede na bajty? Typicky UTF-8.
  \end{description}
  \begin{alertblock}{Přípona nestačí}
    Soubor pojmenovaný \texttt{books.json} nemusí obsahovat platný JSON ani očekávaná data.
  \end{alertblock}
\end{frame}

\begin{frame}{Formát vybíráme podle budoucího čtenáře}
  \begin{enumerate}
    \item Je struktura plochá tabulka, nebo obsahuje vnoření?
    \item Musí data číst člověk, tabulkový procesor, Python, nebo více jazyků?
    \item Potřebujeme zachovat čísla, logické hodnoty a chybějící hodnotu?
    \item Jak velká data budou a musíme je zpracovávat postupně?
    \item Je zdroj důvěryhodný a jak dlouho mají data vydržet?
  \end{enumerate}
  \vspace{2mm}
  \term{Nejdřív požadavky, potom formát.}
\end{frame}

% 0:08--0:19
\begin{frame}{JSON má malý přenositelný datový model}
  \begin{description}
    \item[objekt] dvojice klíč--hodnota; klíč je řetězec,
    \item[pole] uspořádaná posloupnost hodnot,
    \item[skalární hodnoty] řetězec, číslo, \texttt{true}, \texttt{false}, \texttt{null}.
  \end{description}
  \begin{block}{JSON není Python}
    Podobná syntaxe neznamená stejné typy: JSON nezná n-tici, množinu, datum ani vlastní třídu.
  \end{block}
\end{frame}

\begin{frame}{Základní mapování typů}
  \centering
  \begin{tabular}{lll}
    \textbf{Python} & \textbf{JSON} & \textbf{po načtení} \\
    \hline
    \texttt{dict} se stringovými klíči & object & \texttt{dict} \\
    \texttt{list} & array & \texttt{list} \\
    \texttt{tuple} & array & \texttt{list} \\
    \texttt{str} & string & \texttt{str} \\
    \texttt{int}, \texttt{float} & number & číslo \\
    \texttt{True}, \texttt{False} & true, false & \texttt{bool} \\
    \texttt{None} & null & \texttt{None} \\
  \end{tabular}
  \vspace{5mm}

  „Lze zapsat“ neznamená vždy „vrátí se přesně stejný typ“.
\end{frame}

\begin{frame}[fragile]{Ne každý Python objekt projde beze změny}
\begin{lstlisting}[style=pythonSmall]
import json

data = {
    "position": (10, 20),
    "tags": {"python", "data"},
}

text = json.dumps(data)  # TypeError kvuli set
\end{lstlisting}
  \begin{itemize}
    \item N-tice by se po načtení vrátila jako seznam.
    \item Množina není součástí datového modelu JSON.
    \item Datum či vlastní objekt musíme převést podle vlastní dohody.
  \end{itemize}
\end{frame}

\begin{frame}[fragile]{\texttt{json.dump} zapisuje do otevřeného souboru}
\begin{lstlisting}[style=pythonSmall]
import json

with open("books.json", "w", encoding="utf-8") as file:
    json.dump(
        books,
        file,
        ensure_ascii=False,
        indent=2,
    )
\end{lstlisting}
  \begin{itemize}
    \item Režim \texttt{"w"} původní obsah přepíše.
    \item \texttt{json.dump} očekává textový soubor, ne binární.
  \end{itemize}
\end{frame}

\begin{frame}[fragile]{\texttt{json.load} obnoví data ze souboru}
\begin{lstlisting}[style=python]
import json

with open("books.json", encoding="utf-8") as file:
    loaded_books = json.load(file)

print(type(loaded_books))            # list
print(type(loaded_books[0]["year"])) # int
\end{lstlisting}
  Načtený seznam je nový objekt. Změna \texttt{loaded\_books} sama o sobě soubor nezmění.
\end{frame}

\begin{frame}{Dvě souborové a dvě řetězcové funkce}
  \centering
  \begin{tabular}{lll}
    \textbf{Funkce} & \textbf{Vstup nebo výstup} & \textbf{Použití} \\
    \hline
    \texttt{json.dump(data, file)} & otevřený soubor & zápis \\
    \texttt{json.load(file)} & otevřený soubor & čtení \\
    \texttt{json.dumps(data)} & vrací \texttt{str} & text, síť, log \\
    \texttt{json.loads(text)} & přijímá \texttt{str} & text v paměti \\
  \end{tabular}
  \vspace{6mm}

  Písmeno \texttt{s} můžeme číst jako \emph{string}.
\end{frame}

\begin{frame}[fragile]{Čitelnost mění zápis, ne význam dat}
\begin{lstlisting}[style=pythonSmall]
compact = json.dumps(books, ensure_ascii=False)
readable = json.dumps(
    books,
    ensure_ascii=False,
    indent=2,
)
\end{lstlisting}
  \begin{description}
    \item[\texttt{ensure\_ascii=False}] ponechá české znaky čitelné.
    \item[\texttt{indent=2}] přidá odsazení a nové řádky.
  \end{description}
  Čitelnější soubor bývá větší; po načtení však dostaneme stejná data.
\end{frame}

\begin{frame}[fragile]{Poškozený JSON je očekávatelná chyba vstupu}
\begin{lstlisting}[style=pythonSmall]
import json

try:
    with open("books.json", encoding="utf-8") as file:
        books = json.load(file)
except FileNotFoundError:
    print("Soubor neexistuje.")
except json.JSONDecodeError as error:
    print("Neplatny JSON:", error)
\end{lstlisting}
  Zachytáváme konkrétní chyby, kterým rozumíme. Chybu v dalším zpracování neskrýváme širokým \texttt{except Exception}.
\end{frame}

\begin{frame}{Platný JSON ještě nemusí mít správný obsah}
  Oba texty jsou syntakticky platný JSON:
  \begin{columns}[T]
    \begin{column}{.44\textwidth}
      \begin{block}{Očekáváme}
        seznam knih se jmény polí \texttt{title}, \texttt{author}, \texttt{year}
      \end{block}
    \end{column}
    \begin{column}{.44\textwidth}
      \begin{alertblock}{Můžeme dostat}
        číslo \texttt{42}, prázdný objekt nebo knihu bez roku
      \end{alertblock}
    \end{column}
  \end{columns}
  \vspace{5mm}

  Po deserializaci kontrolujeme \term{strukturu, povinná pole, typy a povolené hodnoty}.
\end{frame}

% 0:19--0:31
\begin{frame}{CSV popisuje posloupnost řádků a polí}
  \begin{itemize}
    \item První řádek může obsahovat názvy sloupců.
    \item Každý další řádek představuje jeden záznam.
    \item Oddělovačem nemusí být vždy čárka; běžný je i středník.
    \item Uvozovky dovolují uložit oddělovač nebo nový řádek uvnitř pole.
  \end{itemize}
  \begin{block}{Důležitý důsledek}
    CSV přirozeně vyjadřuje plochou tabulku, ale samo nenese datové typy ani vztahy mezi tabulkami.
  \end{block}
\end{frame}

\begin{frame}[fragile]{\texttt{DictWriter} mapuje slovník na řádek}
\begin{lstlisting}[style=pythonSmall]
import csv

fields = ["title", "author", "year", "available"]

with open("books.csv", "w", encoding="utf-8",
          newline="") as file:
    writer = csv.DictWriter(file, fieldnames=fields)
    writer.writeheader()
    writer.writerows(books)
\end{lstlisting}
  \texttt{fieldnames} určuje sloupce i jejich pořadí.
\end{frame}

\begin{frame}[fragile]{\texttt{DictReader} mapuje řádek na slovník}
\begin{lstlisting}[style=python]
import csv

with open("books.csv", encoding="utf-8",
          newline="") as file:
    reader = csv.DictReader(file)
    loaded_books = list(reader)

print(loaded_books[0])
\end{lstlisting}
  Názvy klíčů převezme z hlavičky. Iterátor lze zpracovávat i postupně bez převodu na seznam.
\end{frame}

\begin{frame}[fragile]{CSV vrací pole jako řetězce}
\begin{lstlisting}[style=python]
book = loaded_books[0]

print(book["year"], type(book["year"]))
# 1965 <class 'str'>

print(book["available"],
      type(book["available"]))
# True <class 'str'>
\end{lstlisting}
  CSV neví, zda \texttt{1965} znamená rok, identifikátor, nebo text. Význam sloupců je součástí naší dohody.
\end{frame}

\begin{frame}[fragile]{Typy obnovujeme explicitně}
\begin{lstlisting}[style=pythonSmall,basicstyle=\ttfamily\small]
def parse_bool(text):
    normalized = text.strip().lower()
    if normalized == "true":
        return True
    if normalized == "false":
        return False
    raise ValueError("Expected true or false")

for book in loaded_books:
    book["year"] = int(book["year"])
    book["available"] = parse_bool(book["available"])
\end{lstlisting}
  Převod je dobré soustředit na jedno místo a při neplatném vstupu selhat srozumitelně.
\end{frame}

\begin{frame}[fragile]{Častá past: \texttt{bool("False")} je \texttt{True}}
\begin{lstlisting}[style=python]
print(bool(""))       # False
print(bool("False"))  # True
print(bool("0"))      # True
\end{lstlisting}
  \begin{block}{Proč?}
    \texttt{bool} u řetězce testuje, zda je neprázdný. Neinterpretuje jeho textový obsah.
  \end{block}
  Pro CSV proto používáme jasně definovaný převod povolených hodnot.
\end{frame}

\begin{frame}[fragile]{CSV neskládáme ručně pomocí \texttt{split}}
\begin{lstlisting}[style=data]
title,author,note
Duna,Frank Herbert,"science fiction, classic"
\end{lstlisting}
  \begin{itemize}
    \item Čárka uvnitř poznámky není nový sloupec, protože pole je v uvozovkách.
    \item Modul \texttt{csv} řeší uvozovky, oddělovače i konce řádků.
    \item Soubor otevíráme s \texttt{newline=""}, aby konce řádků spravoval modul.
  \end{itemize}
  \term{CSV je jednoduché na pohled, ale má pravidla.}
\end{frame}

\begin{frame}{Stejné knihy, jiná dohoda}
  \centering
  \begin{tabular}{lll}
    & \textbf{JSON} & \textbf{CSV} \\
    \hline
    tvar & seznam slovníků & tabulka \\
    čísla a boolean & zachová & načte jako text \\
    vnoření & přirozené & nepohodlné \\
    tabulkový editor & méně pohodlný & přirozený \\
    čtení po řádcích & ne běžný JSON dokument & přirozené \\
    schéma & mimo formát & mimo formát \\
  \end{tabular}
  \vspace{5mm}

  Ani jeden formát není obecně lepší. Každý zpřístupňuje jiné kompromisy.
\end{frame}

% 0:31--0:36
\begin{frame}[fragile]{Vnořená data patří přirozeně do JSONu}
\begin{lstlisting}[style=pythonSmall]
course = {
    "name": "Programovani",
    "students": [
        {"name": "Anna", "score": 92},
        {"name": "Petr", "score": 71},
    ],
}
\end{lstlisting}
  Objekt obsahuje pole dalších objektů. V JSONu zůstane vztah viditelný přímo ve struktuře.
\end{frame}

\begin{frame}{Tabulky vyjadřují vztah pomocí identifikátoru}
  \centering
  \begin{tikzpicture}[>=Latex,node distance=18mm,font=\normalsize,
      box/.style={draw,rounded corners,fill=SoftGray,minimum width=45mm,minimum height=22mm,align=left}]
    \node[box] (machines) {\textbf{machines.csv}\\id, hostname, location\\101, alpha, Lab A};
    \node[box,right=of machines] (users) {\textbf{users.csv}\\id, name, machine\_id\\1, Anna, 101};
    \draw[<->,thick,CodeBlue] (machines)--node[above]{\texttt{machine\_id}}(users);
  \end{tikzpicture}
  \vspace{6mm}

  Vnořenou strukturu lze rozdělit do více tabulek. Vazbu musíme při načtení znovu spojit a ověřit, že odkazované ID existuje.
\end{frame}

% 0:36--0:40
\begin{frame}{Binární formát ukládá bajty, ne čitelný text}
  \begin{itemize}
    \item Může být kompaktní nebo rychlý pro konkrétní účel.
    \item Člověk jej obvykle nezkontroluje v textovém editoru.
    \item Čtenář musí přesně znát formát a jeho verzi.
    \item „Binární“ samo o sobě neznamená menší, rychlejší ani bezpečnější.
  \end{itemize}
  \begin{block}{Příklady}
    Obrázky, komprimované archivy, databázové soubory a Python \texttt{pickle}.
  \end{block}
\end{frame}

\begin{frame}[fragile]{\texttt{pickle} pohodlně uloží Python objekty}
\begin{lstlisting}[style=python]
import pickle

with open("books.pkl", "wb") as file:
    pickle.dump(books, file)

with open("books.pkl", "rb") as file:
    loaded_books = pickle.load(file)
\end{lstlisting}
  Používá binární režimy \texttt{"wb"} a \texttt{"rb"} a je úzce spojený s Pythonem.
\end{frame}

\begin{frame}{Nedůvěryhodný \texttt{pickle} nikdy nenačítáme}
  \begin{alertblock}{Bezpečnostní hranice}
    \texttt{pickle.load} může při načítání spustit libovolný kód. Soubor od cizího člověka, z webu nebo z neověřeného úložiště proto nenačítáme.
  \end{alertblock}
  \begin{itemize}
    \item \texttt{pickle} není bezpečný výměnný formát.
    \item Není vhodný pro komunikaci mezi různými jazyky.
    \item Pro dlouhodobá data může být citlivý na změny tříd a programu.
  \end{itemize}
  Pro výměnu běžných dat preferujeme explicitní formát, například JSON.
\end{frame}

% 0:40--0:43
\begin{frame}{Velikost a rychlost nejdřív měříme}
  \begin{itemize}
    \item Odsazený JSON je větší než kompaktní JSON.
    \item CSV může být úsporné pro plochá data, protože neopakuje názvy polí.
    \item Komprese může rozdíl mezi textovým a binárním formátem zásadně změnit.
    \item Čas často spotřebuje disk, síť nebo další zpracování, ne parser.
  \end{itemize}
  \begin{block}{Praktické pravidlo}
    Nejprve volte podle správnosti, přenositelnosti a jednoduchosti. Optimalizujte až podle měření na skutečných datech.
  \end{block}
\end{frame}

\begin{frame}{Velká data nechceme vždy načíst celá}
  \begin{description}
    \item[CSV] \texttt{DictReader} lze procházet řádek po řádku.
    \item[JSON] Jeden velký objekt či pole \texttt{json.load} běžně načte do paměti celé.
    \item[JSON Lines] Každý řádek je samostatný JSON objekt; lze jej zpracovat postupně.
  \end{description}
  \begin{alertblock}{Pozor}
    JSON Lines je jiná dohoda než jeden běžný JSON dokument. Odesílatel i příjemce ji musí znát.
  \end{alertblock}
\end{frame}

\begin{frame}{Rozhodovací mapa}
  \centering
  \begin{tabular}{ll}
    \textbf{Potřeba} & \textbf{Rozumný výchozí formát} \\
    \hline
    jednoduchá tabulka & CSV \\
    otevření v tabulkovém procesoru & CSV \\
    vnořená struktura a běžné typy & JSON \\
    výměna mezi různými jazyky & JSON \\
    postupné zpracování záznamů & CSV nebo JSON Lines \\
    vlastní důvěryhodný Python mezivýsledek & někdy \texttt{pickle} \\
    dlouhodobá či důležitá data & dokumentovaný stabilní formát \\
  \end{tabular}
\end{frame}

% 0:43--0:45
\begin{frame}{Časté chyby}
  \begin{itemize}
    \item Přepsání jediného zdrojového souboru režimem \texttt{"w"} před ověřením výsledku.
    \item Záměna \texttt{dump}/\texttt{dumps} nebo \texttt{load}/\texttt{loads}.
    \item Předpoklad, že CSV obnoví čísla a logické hodnoty samo.
    \item Použití \texttt{bool(text)} pro převod textové logické hodnoty.
    \item Ruční rozdělování CSV pomocí \texttt{split(",")}.
    \item Důvěra v příponu souboru bez ověření obsahu.
    \item Načtení nedůvěryhodného souboru pomocí \texttt{pickle}.
  \end{itemize}
\end{frame}

\begin{frame}{Rychlá kontrola porozumění}
  \begin{enumerate}
    \item Jaký je rozdíl mezi serializací a zápisem souboru?
    \item Co vrátí JSON po uložení a načtení n-tice?
    \item Proč je rok načtený přes \texttt{DictReader} řetězec?
    \item Proč \texttt{bool("False")} nepřevede text podle očekávání?
    \item Kdy je přirozenější CSV a kdy JSON?
    \item Za jakých podmínek lze bezpečně použít \texttt{pickle.load}?
  \end{enumerate}
\end{frame}

\begin{frame}{Shrnutí}
  \begin{itemize}
    \item Serializace převádí datový model na text nebo bajty; deserializace vytváří nové objekty.
    \item JSON zachovává běžné typy a přirozeně vyjadřuje vnoření.
    \item CSV je přirozená tabulka, ale typy sloupců obnovujeme sami.
    \item Formát volíme podle tvaru dat, příjemce, velikosti a důvěryhodnosti.
    \item Platná syntaxe nezaručuje správný obsah; načtená data validujeme.
    \item Nedůvěryhodný \texttt{pickle} nikdy nenačítáme.
  \end{itemize}
  \begin{block}{Navazující komentované příklady}
    Porovnáme JSON a CSV nad stejnými daty, obnovíme typy a převedeme tabulku do vnořené struktury.
  \end{block}
\end{frame}

\end{document}
